\documentclass[10pt]{article}
\usepackage[T1]{fontenc}
\usepackage[utf8]{inputenc}
\usepackage[french]{babel}
\usepackage{amsmath,amssymb,mathrsfs,mathpazo}
\usepackage{geometry}
\geometry{a4paper,margin=2cm}
\usepackage{graphicx,array,tabularx,tabularray,colortbl,multirow,multicol,booktabs,enumitem,pifont,ulem,fancybox,diagbox,cancel}
\usepackage{pstricks,pstricks-add,pst-plot,pst-tree,pst-node,pst-text,pst-func,pst-math,pst-bspline,pst-3dplot,pst-eucl,pst-solides3d}
\let\pscalculate\undefined
\usepackage{pst-bezier}
\usepackage{tikz,pgfplots,tkz-tab}
\usetikzlibrary{arrows,arrows.meta,patterns,calc,positioning,trees,decorations.markings,shapes.geometric,shapes.misc,shapes.arrows,intersections,decorations.pathmorphing,decorations.pathreplacing,backgrounds,shadows,babel}
\pgfplotsset{compat=1.18}
\usepgfplotslibrary{fillbetween}
\usepackage[np]{numprint}
\usepackage{listings,makecell,pas-tableur,twoopt}
\let\barre\undefined
\usepackage[pstricks]{tablvar}
\definecolor{darkgreen}{RGB}{0,100,0}
\providecommand{\red}{\color{red}}
\providecommand{\blue}{\color{blue}}
\providecommand{\green}{\color{green}}
\providecommand{\black}{\color{black}}
\providecommand{\white}{\color{white}}
\providecommand{\gray}{\color{gray}}
\providecommand{\cyan}{\color{cyan}}
\providecommand{\magenta}{\color{magenta}}
\providecommand{\yellow}{\color{yellow}}
\providecommand{\violet}{\color{violet}}
\providecommand{\orange}{\color{orange}}
\providecommand{\euro}{\texteuro}
\providecommand{\EUR}{\texteuro}
\providecommand{\ieme}{\textsuperscript{e}}
\providecommand{\degres}{\ensuremath{^\circ}}
\providecommand{\degree}{\ensuremath{^\circ}}
\providecommand{\vv}[1]{\overrightarrow{#1}}
\providecommand{\Coord}[1]{#1}
\providecommand{\decosix}{}
\providecommand{\starredbullet}{\ensuremath{\bullet}}
\pagestyle{empty}
\setlength{\parindent}{0pt}
\usepackage{adjustbox}\DeclareUnicodeCharacter{2236}{:}\DeclareUnicodeCharacter{2190}{\ensuremath{\leftarrow}}
\providecommand{\R}{\mathbb{R}}
\providecommand{\N}{\mathbb{N}}
\providecommand{\Z}{\mathbb{Z}}
\providecommand{\Q}{\mathbb{Q}}
\providecommand{\D}{\mathbb{D}}
\providecommand{\C}{\mathbb{C}}
\providecommand{\e}{{\mathrm{e}}}
\providecommand{\vect}[1]{\overrightarrow{#1}}
\providecommand{\vectt}[1]{\overrightarrow{\mathrm{#1}}}
\providecommand{\vvt}[1]{\overrightarrow{\mathrm{#1}}}
\providecommand{\np}[2][]{{#2~#1}}
\providecommand{\vv}[1]{{\overrightarrow{#1}}}
\renewcommand{\d}{{\mathrm{d}}}
\providecommand{\degres}{{^\circ}}
\providecommand{\degre}{{^\circ}}
\providecommand{\blue}{\color{blue}}
\providecommand{\red}{\color{red}}
\providecommand{\enskip}{\quad}
\providecommand{\llbracket}{\lbrack\!\lbrack}
\providecommand{\rrbracket}{\rbrack\!\rbrack}
\providecommand{\euro}{{\text{€}}}
\providecommand{\no}{{\text{n°}}}
\providecommand{\interval}[2]{{[#1;#2]}}
\providecommand{\intervFO}[2]{{[#1;#2[}}
\providecommand{\dotfill}{\ldots}

\providecommand{\ds}{\displaystyle}
\providecommand{\textding}[1]{#1}
\providecommand{\Oij}{\(\mathrm{O};\vec{\imath},\vec{\jmath}\)}
\providecommand{\Oijk}{\(\mathrm{O};\vec{\imath},\vec{\jmath},\vec{k}\)}
\providecommand{\Ouv}{\(\mathrm{O};\vec{u},\vec{v}\)}
\providecommand{\intervFF}[2]{\ensuremath{[#1;#2]}}
\renewcommand{\euro}{\ifmmode\text{\texteuro}\else\texteuro\fi}
\DeclareUnicodeCharacter{00B0}{\ifmmode^{\circ}\else\textdegree\fi}
\DeclareUnicodeCharacter{20AC}{\ifmmode\text{\texteuro}\else\texteuro\fi}
% Énoncés : Éducation nationale. Transcription et corrigés : APMEP, Denis Vergès et contributeurs. Découpage et classement : CoopMaths, Sébastien Lozano, Éric Elter et contributeurs. Adaptation de la présentation : Maths Explorer.
% CC BY-SA : https://coopmaths.fr/www/about/
% https://coopmaths.fr/alea/static/bac/2021/tex/bac_2021_06_polynesie_4.tex
\begin{document}
\noindent\textit{APMEP / CoopMaths. CC BY-SA. Présentation adaptée par Maths Explorer.}\par\medskip
\medskip

\begin{tabularx}{\linewidth}{|X|}\hline
\textbf{Principaux domaines abordés : Fonction exponentielle, convexité, dérivation, }\\ \textbf{équations différentielles}\\ \hline
\end{tabularx}

\medskip

Cet exercice est composé de trois parties indépendantes.

\medskip

On a représenté ci-dessous, dans un repère orthonormé, une portion de la courbe
représentative $\mathcal{C}$ d'une fonction $f$ définie sur $\R$ :

\begin{center}
\psset{unit=1.25cm}
\begin{pspicture*}(-3,-0.95)(3.5,3.5)
\psgrid[gridlabels=0pt,subgriddiv=1,griddots=8]
\psaxes[linewidth=1.25pt,labelFontSize=\scriptstyle]{->}(0,0)(-3,-0.95)(3.5,3.5)
\psplot[plotpoints=2000,linewidth=1.25pt,linecolor=red]{-2.25}{3.5}{x 2 add 2.71828 x exp div}
\psplot[plotpoints=200,linewidth=1.5pt,linestyle=dotted]{-1.5}{2.5}{ 2 x sub}
\uput[ul](-1.6,2){\red $\mathcal{C}$}\uput[ur](0,2){A}\uput[ur](2,0){B}
\uput[d](3.4,0){\small $x$}\uput[l](0,3.4){\small $y$}
\end{pspicture*}
\end{center}

On considère les points A(0~;~2) et B(2~;~0).

\bigskip

\begin{center}\textbf{Partie 1}\end{center}

Sachant que la courbe $\mathcal{C}$ passe par A et que la droite (AB) est la tangente à la courbe $\mathcal{C}$ au point A, donner par lecture graphique : 

\medskip

\begin{enumerate}
\item La valeur de $f(0)$ et celle de $f'(0)$.
\item Un intervalle sur lequel la fonction $f$ semble convexe.
\end{enumerate}

\bigskip

\begin{center}\textbf{Partie 2}\end{center}

On note $(E)$ l'équation différentielle 

\[y' = -y + \text{e}^{-x}.\]

On admet que $g :\,  x \longmapsto  x\text{e}^{-x}$ est une solution particulière de $(E)$.

\medskip

\begin{enumerate}
\item Donner toutes les solutions sur $\R$ de l'équation différentielle $(H)$ : 
$y' = -y$.
\item En déduire toutes les solutions sur $\R$ de l'équation différentielle $(E)$.
\item Sachant que la fonction $f$ est la solution particulière de $(E)$ qui vérifie $f(0) = 2$, déterminer une expression de $f(x)$ en fonction de $x$.
\end{enumerate}

\bigskip

\begin{center}\textbf{Partie 3}\end{center}

On admet que pour tout nombre réel $x$,\, $f(x) = (x + 2)\text{e}^{-x}$.

\medskip

\begin{enumerate}
\item On rappelle que $f'$ désigne la fonction dérivée de la fonction $f$.
	\begin{enumerate}
		\item Montrer que pour tout $x \in \R,\, f'(x) = (-x - 1) \text{e}^{-x}$.
		\item Étudier le signe de $f'(x)$ pour tout $x \in \R$ et dresser le tableau des variations de $f$ sur $\R$.

On ne précisera ni la limite de $f$ en $- \infty$ ni la limite de $f$ en $+ \infty$.

On calculera la valeur exacte de l'extremum de $f$ sur $\R$.
	\end{enumerate}
\item On rappelle que $f''$ désigne la fonction dérivée seconde de la fonction $f$.
	\begin{enumerate}
		\item Calculer pour tout $x \in \R,\, f''(x)$.
		\item Peut-on affirmer que $f$ est convexe sur l'intervalle $[0~;~+\infty[$?
	\end{enumerate}
\end{enumerate}

\bigskip

\end{document}
